\documentclass[../main.tex]{subfiles}
\begin{document}
Touchscreen interaction is widely used in modern mobile devices, yet it is typically limited to two-dimensional touch inputs, restricting richer interaction capabilities. Recent studies have explored the use of capacitive images to identify which finger interacts with the screen, with the CapFingerId dataset providing a large-scale benchmark for this task. While convolutional neural networks (CNNs) achieve high classification accuracy, they often require significant computational resources. This thesis therefore investigates energy-efficient alternatives using Spiking Neural Networks (SNNs), including a model developed on the RANC\cite{RANC} framework with potential for FPGA deployment, although hardware implementation is beyond the scope of this work. In addition, a SpikeFormer-based SNN Vision Transformer is explored, leveraging direct training and strong representation learning capabilities instead of ANN-to-SNN conversion. This chapter presents the problem context, reviews related works, and introduces the necessary theoretical background.

\section{Scope of Research}
\subsection{Task Description}

This thesis addresses the problem of classifying capacitive fingerprint images collected from touchscreen devices. Unlike traditional input modalities that only capture 2D touch locations, capacitive sensing enables the acquisition of low-resolution fingerprint-like images, which contain distinctive patterns related to different fingers. The CapFingerId \cite{CapfingerId} dataset provides a large-scale collection of such images, comprising various finger combinations, particularly focusing on commonly used fingers such as the thumb and index finger.

However, capacitive images present several challenges compared to conventional optical fingerprint images. Specifically, they are characterized by low resolution, high noise levels, and significant variations in touch position and pressure. Furthermore, the task requires the model to learn position-invariant features in order to robustly distinguish between different fingers under varying interaction conditions. These characteristics make the classification problem non-trivial and motivate the need for robust and efficient learning approaches.

In this thesis, both the proposed RANC-based Spiking Neural Network (SNN) and the Spiking Vision Transformer (Spiking ViT) are designed to perform multi-class finger classification on capacitive touch images. The classification tasks include distinguishing between:

\begin{enumerate}[label=\Roman*.]
    \item Left and right thumbs,
    \item Left and right index fingers,
    \item Thumb and index finger categories,
    \item Thumb and non-thumb finger categories,
    \item Different fingers belonging to the same hand.
\end{enumerate}

These classification settings enable comprehensive evaluation of the proposed architectures under multiple levels of finger discrimination difficulty, ranging from coarse-grained categorization to fine-grained finger identification.

Although the CapFingerId dataset was collected under multiple interaction activities such as tapping, swiping, and other touchscreen gestures, this research focuses exclusively on samples generated from \textit{tap} interactions. The tapping scenario provides relatively stable finger contact patterns and allows more controlled evaluation of the proposed models without the additional motion variability introduced by dynamic gestures.

\subsection{Problem Formulation}
Formally, let the dataset be defined as:
\begin{equation}
D = \{(x_i, y_i)\}_{i=1}^{N}
\end{equation}
where $x_i$ denotes the input capacitive image and $y_i \in \{1, \ldots, K\}$ represents the corresponding class label. The objective is to learn a mapping function:
\begin{equation}
f(x; \theta) \rightarrow y
\end{equation}
where $\theta$ denotes the model parameters. The model is trained to minimize the cross-entropy loss for classification.

\subsection{Evaluation Metrics}

The performance of the proposed models is evaluated primarily using classification accuracy. Accuracy measures the proportion of correctly classified samples over the total number of evaluated samples and is defined as:

\[
\text{Accuracy} =
\frac{N_{\text{correct}}}{N_{\text{total}}}
\]

where \(N_{\text{correct}}\) denotes the number of correctly predicted samples and \(N_{\text{total}}\) represents the total number of testing samples.

Accuracy is selected as the primary evaluation metric because the classification tasks considered in this thesis are formulated as multi-class finger classification problems with balanced evaluation settings. The metric provides a straightforward and interpretable measure of the overall recognition capability of the proposed architectures.

In addition to classification accuracy, the number of trainable parameters is also reported to evaluate the computational complexity of different model configurations. Since the proposed architectures are intended for lightweight and resource-constrained neuromorphic systems, parameter efficiency is considered an important factor alongside recognition performance.

Therefore, the experimental analysis focuses on the trade-off between classification accuracy and model complexity under different architectural and temporal encoding configurations.

\subsection{Constraints and Research Gap}
While convolutional neural networks (CNNs) have demonstrated strong performance on similar tasks, they often require significant computational resources. In contrast, Spiking Neural Networks (SNNs) offer a more energy-efficient alternative but remain underexplored for capacitive fingerprint classification. Furthermore, the application of transformer-based models within the SNN paradigm is still limited. These challenges motivate the development of efficient and accurate models in this thesis.

\section{Related Works}
\subsection{Fingerprint Recognition using Capacitive Data}
\subsubsection{Motivation}
Early approaches to finger identification often relied on external sensors or additional hardware, which introduces practical limitations in terms of cost, usability, and deployment \cite{CapfingerId}. In contrast, capacitive touchscreen technology inherently captures variations in electrical capacitance caused by the ridge and valley structures of human fingerprints, allowing it to generate fingerprint-like patterns without requiring dedicated biometric sensors \cite{review_capacitive, A_Review_of_Fingerprint_Sensors}. This capability has motivated recent research to utilize capacitive images as a lightweight and readily available source of biometric information for finger classification and user interaction.

Furthermore, prior work has demonstrated that capacitive images can be effectively leveraged using deep learning techniques. For instance, the CapFingerId study shows that convolutional neural networks can achieve over 92\% accuracy in distinguishing between different fingers using only low-resolution capacitive images \cite{CapfingerId}. Similarly, more recent work such as SpeciFingers highlights the potential of learning spatio-temporal features from capacitive data to improve finger identification performance \cite{SpeciFingers}. These results indicate that capacitive fingerprint recognition is both feasible and promising for enhancing human-computer interaction.

Despite these advances, capacitive fingerprint data remains challenging due to its low resolution, high noise levels, and sensitivity to touch variations such as position and pressure. These challenges motivate the exploration of more robust and efficient models, particularly energy-efficient approaches such as Spiking Neural Networks, which are better suited for deployment on edge devices.

\subsubsection{CapFingerId Dataset}

The CapFingerId dataset~\cite{CapfingerId} was introduced as one of the first publicly available datasets for finger recognition using capacitive images captured from commodity smartphone touchscreens. The dataset was released together with baseline CNN models to facilitate research on learning-based finger identification from low-resolution capacitive sensing data.

The dataset was collected from 20 participants (15 male and 5 female) with a mean age of 22.4 ($Std = 2.8$). Among them, 18 were right-handed and 2 were left-handed. Hand length ranged from 16.3~cm to 20.8~cm ($M = 18.9$~cm, $Std = 1.3$). Data acquisition was performed on an \textbf{LG Nexus~5} smartphone equipped with a \textbf{Synaptics ClearPad~3350} mutual-capacitive touchscreen. The modified firmware enabled access to raw capacitive images with a sensor resolution of $15 \times 27$, where each sensing cell corresponds to approximately $4.1 \times 4.1$~mm on the touchscreen. Capacitive images were sampled at 20 frames per second during user interactions.

Participants performed three representative touch interactions using all ten fingers: tapping, dragging, and scrolling. Each finger completed 30 repetitions of every task, with the scrolling task equally divided into horizontal and vertical directions, resulting in 900 trials per participant. Touch locations were randomized, and finger order was balanced using a Latin square to reduce ordering effects. Capacitive images were automatically labeled according to the instructed finger and interaction task.

After removing empty samples and frames containing multiple simultaneous touches, blob detection was applied to isolate individual contact regions. Each detected blob was cropped and re-centered into a $15 \times 27$ matrix, producing a total of 455,709 labeled capacitive images. Unlike sequence-based datasets, each sample consists of a single capacitive frame representing one finger contact.

The dataset supports multiple finger classification scenarios, including left-versus-right thumb, left-versus-right index finger, thumb-versus-index, thumb-versus-other fingers, five-finger classification, and full ten-finger classification. Using convolutional neural networks as baseline models, the authors demonstrated that reliable recognition of all ten fingers remains challenging because of the limited spatial resolution of smartphone capacitive sensors. Nevertheless, distinguishing between the left and right thumbs achieved an accuracy exceeding 92\%, establishing CapFingerId as an important benchmark for capacitive finger recognition on commodity mobile devices.

\begin{figure}[h]
    \centering
    \begin{subfigure}{0.3\textwidth}
        \centering
        \includegraphics[width=\textwidth]{Figure/hand.jpg}
        \caption{Using two thumbs for interaction}
    \end{subfigure}
    \hspace{1cm}
    \begin{subfigure}{0.15\textwidth}
        \centering
        \includegraphics[width=\textwidth]{Figure/cap_img.png}
        \caption{Raw Data}
    \end{subfigure}
    \caption{ Identifying left and right thumbs on a commoditysmartphone using the raw capacitive data of touches.}
    \label{fig:two_images}
\end{figure}

\begin{figure}[h]
    \centering
    \includegraphics[width=0.6\textwidth]{Figure/CapfingerId_video.png}
    \caption{Volunteers carried out predefined tasks for data collection purposes}
    \label{fig:my_image}
\end{figure}

\begin{figure}[t]
    \centering
    % Hàng 1
    \begin{minipage}{0.15\linewidth}
        \centering
        
\includegraphics[width=\linewidth]{Figure/0.png}
        \label{fig:sample1}
    \end{minipage}
    \hspace{0.02\linewidth}
    \begin{minipage}{0.15\linewidth}
        \centering
        
\includegraphics[width=\linewidth]{Figure/1.png}
        \label{fig:sample2}
    \end{minipage}

    \vspace{0.5em} % khoảng cách giữa 2 hàng

    % Hàng 2
    \begin{minipage}{0.15\linewidth}
        \centering
        
\includegraphics[width=\linewidth]{Figure/2.png}
        \label{fig:sample3}
    \end{minipage}
    \hspace{0.02\linewidth}
    \begin{minipage}{0.15\linewidth}
        \centering
        
\includegraphics[width=\linewidth]{Figure/3.png}
        \label{fig:sample4}
    \end{minipage}

    \caption{Representative samples from dataset.}
    \label{fig:dataset_samples}
\end{figure}

\subsubsection{SpeciFingers Dataset}
To support the training and evaluation of finger identification models on tabletop 
capacitive touchscreens, the SpeciFingers~\cite{huang2024specifingers} dataset was 
constructed and publicly released alongside the corresponding model and source code.

The data were collected from 20 participants (13 male and 7 female) with a mean age 
of 27.6 ($Std = 4.60$). Among the participants, 18 were right-handed and 2 were 
left-handed. Hand length, measured from the wrist crease to the tip of the middle 
finger, ranged from 15.8~cm to 19.8~cm ($M = 18.37$~cm, $Std = 1.12$). The 
experimental device was a \textbf{Wacom DTH-2242} tabletop touchscreen with a 21.5-inch 
display ($475.2$~mm $\times$ $267.3$~mm) and a screen resolution of 
$1920 \times 1080$ pixels. Raw touch sensitivity data was accessed via \textbf{Wacom's 
Feel\texttrademark{} Multi-Touch SDK} and recorded at a rate of 100 frames per 
second (every 10~ms). The touchscreen provides a capacitive sensor resolution of 
$122 \times 70$, with sensitivity values ranging from 0 to 65,535. The device was 
placed on a desk at an angle of $20^\circ$, and participants were instructed to 
hover their hands naturally over the screen during interactions.

Participants performed four types of single-finger touch interactions using all ten 
fingers: tapping, pressing, sliding (in four directions: up, down, left, and right), 
and dragging. The tapping task required participants to tap a $50 \times 50$~px 
target appearing at randomized positions. The pressing task required sustaining 
contact for at least 0.5 seconds. The sliding task covered four cardinal directions, 
and the dragging task involved moving a target element to a destination position. 
Each task was repeated 40 times per finger, resulting in $10~\text{fingers} \times 
4~\text{tasks} \times 40~\text{repetitions} = 1{,}600$ trials per participant. The 
order of fingers was balanced using a \textbf{Balanced Latin Square}, and task order was 
randomized for each finger. Recording of each action was triggered automatically 
upon finger contact and terminated upon finger release, eliminating the need for 
manual segmentation.

After filtering incorrectly labeled trials and removing sequences containing 
unintentional touches via blob detection, the final dataset comprises 38,153 
sequences of capacitive images. Each raw sensitivity frame is represented as a 
$122 \times 70$ matrix, from which contact blobs are extracted and resized to 
$224 \times 224$ pixels for use as model input. The average sequence length varies 
by action type: approximately 10.81 frames for tapping, 20.46 frames for sliding, 
74.50 frames for dragging, and 84.26 frames for pressing.

The dataset is labeled for all ten fingers across four action types, supporting a 
range of classification scenarios. The authors explored three finger grouping 
configurations: classification across all ten fingers of both hands 
($\text{ALL}_{L\&R}$), all five fingers of a single hand ($\text{ALL}_{L/R}$), 
and a three-class grouping consisting of the thumb, little finger, and the remaining 
three middle fingers combined into a single category ($TM^{+}L$). The last 
configuration was motivated by the observation that the index, middle, and ring 
fingers produce highly similar capacitive patterns that are difficult to distinguish 
reliably.

Compared to the previously available \textbf{CapFingerId} dataset~\cite{CapfingerId}, 
\textbf{SpeciFingers} is larger in sequence count (38,153 vs.\ 17,416), covers a greater 
variety of action types (4 vs.\ 3), and was collected on a higher-resolution 
tabletop device rather than a smartphone, making it particularly relevant for 
research on finger-specific interaction in tabletop computing environments.

\subsubsection{Learning-Based Finger Recognition}

Following the release of the CapFingerId dataset, Le \textit{et al.}~\cite{CapfingerId} established one of the first deep learning baselines for capacitive image-based finger recognition. Prior to training deep models, the authors explored handcrafted geometric descriptors extracted from capacitive touch blobs, including blob area, average capacitance, ellipse dimensions, and orientation. These features were evaluated using conventional machine learning classifiers such as Random Forests but exhibited limited discriminative capability, particularly for distinguishing fingers with highly similar touch patterns. Motivated by these limitations, the authors proposed a convolutional neural network (CNN) that directly learns spatial representations from raw capacitive images, eliminating the need for manual feature engineering.

The proposed CNN accepts a single $15 \times 27$ capacitive image as input and predicts the corresponding finger class. The model was evaluated under multiple classification scenarios, including left-versus-right thumb, left-versus-right index finger, thumb-versus-index, thumb-versus-other fingers, five-finger, and ten-finger classification. Experimental results demonstrated that CNN-based feature learning consistently outperformed handcrafted feature-based approaches across all tasks. In particular, the model achieved over 92% accuracy for distinguishing between the left and right thumbs, while full ten-finger classification remained challenging because the low spatial resolution of commodity smartphone capacitive sensors provides insufficient discriminative information. These findings established CNNs as an effective baseline for capacitive finger recognition and highlighted the need for more powerful architectures capable of extracting richer representations from low-resolution capacitive images.

Following the CapFingerId study, Huang et al.~\cite{huang2024specifingers}
argued that finger touch data should be regarded as a temporal sequence rather
than a collection of independent capacitive images. They observed that the
capacitive response evolves throughout an interaction, with the contact area
and sensitivity values changing continuously over time. Based on this
observation, they proposed an encoder--decoder architecture that jointly
models spatial and temporal information from capacitive image sequences
instead of performing frame-wise classification. During preprocessing,
contact blobs were extracted from each frame, resized to a fixed resolution,
and normalized to a uniform sequence length through interpolation or temporal
sampling. Online data augmentation, including random flipping and rotation,
was also applied to improve model generalization.

The proposed architecture consists of a CNN encoder followed by an LSTM
decoder. The CNN encoder extracts spatial representations from individual
capacitive images, while the LSTM aggregates temporal dependencies across
consecutive frames to produce the final finger prediction. Experimental
results on both the SpeciFingers and CapFingerId datasets demonstrated that
incorporating temporal information consistently improves recognition
performance over the frame-based CNN baseline. In particular, the proposed
model improved the three-class ($TM^{+}L$) classification accuracy by nearly
10 percentage points on the SpeciFingers dataset and by more than 12
percentage points on the CapFingerId dataset. These results demonstrate that
modeling the temporal dynamics of capacitive image sequences is an effective
strategy for capacitive image-based finger recognition.

In this thesis, the CNN model proposed by Le \textit{et al.}~\cite{CapfingerId}
is adopted as the primary baseline for evaluating the proposed models.
This choice is motivated by the fact that it represents the first deep learning
baseline for capacitive image-based finger recognition and has been extensively
evaluated on the CapFingerId dataset, which is also used throughout this work.
Its relatively simple architecture provides a fair and well-established
reference, allowing the effectiveness of the proposed models to be assessed
through direct comparison under the same experimental setting.

\subsection{Spiking Neural Networks (SNN)}

Spiking Neural Networks (SNNs) represent the third generation of neural network models, inspired by the dynamics of biological neurons. Unlike conventional artificial neural networks (ANNs) that process continuous values, SNNs operate using discrete spike events over time. This event-driven mechanism allows computations to occur only when spikes are present, leading to significantly reduced energy consumption. As a result, SNNs have emerged as a promising approach for energy-efficient computing, particularly in edge devices and neuromorphic hardware systems.

\subsubsection{Neuron Model}

A fundamental component of SNNs is the spiking neuron model. In this work, the Leaky Integrate-and-Fire (LIF) neuron is adopted due to its simplicity and effectiveness. The membrane potential of a neuron accumulates incoming spikes over time and decays gradually. When the membrane potential exceeds a predefined threshold, the neuron emits a spike and resets its state. This behavior can be described as:

\begin{equation}
\tau \frac{dV(t)}{dt} = -V(t) + I(t)
\end{equation}

where $V(t)$ is the membrane potential, $I(t)$ is the input current, and $\tau$ is the membrane time constant.

\subsubsection{Spike Encoding}

To process static image data within the SNN framework, it is necessary to convert input images into spike trains. In this work, \textit{burst coding} is adopted as the primary encoding scheme. Unlike traditional rate coding, burst coding represents input intensity using short bursts of spikes, where both the number of spikes within a burst and their temporal spacing carry information. Stronger input signals typically generate bursts with higher spike counts and shorter inter-spike intervals. This encoding strategy enables richer temporal dynamics while maintaining relatively low latency.

Burst coding offers several advantages over conventional rate coding. By transmitting multiple spikes in rapid succession, it can improve robustness to noise and enhance signal reliability, especially in early layers of the network \cite{panda2016unsupervisedregenerativelearninghierarchical}. Moreover, burst-based representations can achieve a better trade-off between energy efficiency and information fidelity, as meaningful information is conveyed in compact temporal patterns rather than long spike trains \cite{park2019fastefficientinformationtransmission}.

In addition to burst coding, several alternative encoding schemes have been proposed to improve efficiency and temporal representation. \textit{Rate coding} encodes pixel intensity as the firing frequency over a fixed time window and is widely used due to its simplicity and robustness \cite{Gerstner_Kistler_2002, diehl2015unsupervised}. However, it often requires longer simulation time to achieve high accuracy.

\textit{Temporal coding} encodes information in the precise timing of spikes rather than their frequency. For example, in \textit{time-to-first-spike} coding, stronger input signals generate earlier spikes, allowing for faster inference with fewer spikes \cite{thorpe2001spike, masquelier2007unsupervised}. This approach is more biologically plausible and can significantly reduce latency.

Another related approach is \textit{latency coding}, where input intensity is represented by the delay before a neuron emits its first spike. While effective for low-latency systems, it may be sensitive to input noise \cite{bohte2002error}.

\textit{Population coding} represents a single input value using a group of neurons with overlapping receptive fields, enabling a distributed and potentially more robust representation \cite{eliassmith2003neural}. However, this comes at the cost of increased neuron count and computational complexity.

Finally, \textit{event-based encoding} is commonly used in neuromorphic vision sensors, where spikes are generated only when changes in the input signal occur. Although highly efficient, it is less directly applicable to static image datasets such as CapFingerId \cite{gallego2020event}.

Overall, burst coding is selected in this work due to its ability to capture both rate and temporal information, providing a balanced trade-off between accuracy, latency, and energy efficiency for static image-based SNN applications.

\subsubsection{Comparison of Spiking Neuron Models}

Various spiking neuron models have been proposed in the literature, each offering different trade-offs between biological realism, computational efficiency, and suitability for deep learning.

The \textbf{Integrate-and-Fire (IF) model} is the simplest formulation, in which the membrane potential accumulates incoming inputs without any leakage mechanism. Its main advantage lies in its computational efficiency and ease of implementation. However, the absence of a decay term makes it less biologically plausible and can lead to unbounded accumulation of membrane potential, which may negatively affect network stability \cite{Gerstner_Kistler_2002}.

The Leaky \textbf{Integrate-and-Fire (LIF) model} extends IF by incorporating a leakage term that gradually decays the membrane potential over time. This simple yet effective modification significantly improves temporal dynamics while maintaining low computational cost. As a result, LIF achieves a favorable balance between biological plausibility and efficiency, making it one of the most widely adopted neuron models in modern SNN research \cite{Gerstner_Kistler_2002, neftci2019surrogate}. Furthermore, LIF neurons integrate naturally with surrogate gradient methods, enabling stable end-to-end training in deep spiking networks \cite{neftci2019surrogate}.

More advanced variants such as the \textbf{Parametric LIF (PLIF)} introduce learnable parameters, such as the membrane time constant, allowing the neuron to adapt its temporal dynamics during training. This flexibility can improve performance across different tasks and network depths, but comes at the cost of increased model complexity and additional training overhead \cite{fang2021incorporating}.

Similarly, the \textbf{Adaptive LIF (ALIF) model} incorporates dynamic threshold modulation to simulate neuronal adaptation. This enables better modeling of long-term temporal dependencies, but also increases computational complexity and may complicate optimization \cite{bellec2018long}.

At the other end of the spectrum, biologically detailed models such as the \textbf{Izhikevich and Hodgkin--Huxley neurons} provide highly accurate representations of neuronal behavior. These models are capable of reproducing a wide range of spiking patterns observed in biological systems. However, their high computational cost and complex dynamics make them impractical for large-scale deep learning applications, particularly in vision-based tasks \cite{izhikevich2003simple, hodgkin1952quantitative}.

Considering the trade-offs discussed above, the \textbf{Leaky Integrate-and-Fire (LIF) neuron} is adopted in this work due to its balance between computational efficiency and temporal modeling capability. Its simple formulation enables efficient implementation and stable training using surrogate gradient methods, which is crucial for large-scale deep SNN architectures. While more advanced neuron models offer greater expressiveness, they often introduce unnecessary complexity for the scope of this study. Therefore, LIF is selected to ensure scalability and efficiency, with further exploration of more sophisticated neuron dynamics left for future work.

\subsubsection{Advantages and Challenges}

SNNs offer several advantages over traditional neural networks. Their event-driven nature enables sparse computation, resulting in lower power consumption and improved efficiency, especially when deployed on neuromorphic hardware. However, SNNs also present several challenges, including difficulties in training, sensitivity to encoding schemes, and potential performance gaps compared to conventional CNNs.

Despite these challenges, the potential of SNNs for energy-efficient learning makes them a compelling choice for this work. In particular, their compatibility with neuromorphic frameworks motivates the exploration of hardware-aware implementations, as discussed in the following section.
\subsection{Neuromorphic Hardware \& Framework (RANC)}
\subsubsection{Overview of RANC\cite{RANC} Architecture}
The \textbf{Reconfigurable Architecture for Neuromorphic Computing (RANC)} is a highly scalable and flexible hardware ecosystem designed to optimize the execution of Spiking Neural Networks (SNNs). Unlike traditional von Neumann architectures, RANC\cite{RANC} adopts a distributed and event-driven design paradigm, where computations are performed locally within processing units and communication occurs only when spike events are generated.

From a structural perspective, the RANC\cite{RANC} system is composed of multiple interconnected processing elements, referred to as neuron cores. Each core integrates a group of spiking neurons along with their associated synaptic weights and local SRAM memory for storing configuration parameters. These cores operate in parallel and are synchronized through a global clock signal (global tick), enabling efficient large-scale neural computation.

\begin{figure}[tb]
    \centering
    \includegraphics[width=0.95\linewidth]{Figure/ComponentDiagram.pdf}
    \caption{High level architectural overview of RANC\cite{RANC} components.}
    \vspace{-12pt}
    \label{fig:arch_component_overview}
\end{figure}

Communication between neuron cores is implemented through an on-chip network (NoC) based on a 2D mesh topology Figure~\ref{fig:arch_component_overview}, utilizing a packet-based routing mechanism. In this scheme, when a neuron generates a spike, the destination information is retrieved from the core’s SRAM to construct a packet. Each packet contains four key fields: signed relative offsets $dx$ and $dy$ used for routing the packet to the target core, an axon index to identify the receiving neuron, and a delivery tick offset that specifies the time at which the spike will be processed by the scheduler. The event-driven communication mechanism, combined with deterministic XY routing, minimizes unnecessary data movement and significantly improves the overall energy efficiency of the system.

A key advantage of RANC\cite{RANC} lies in its strong reconfigurability. The architecture allows flexible mapping of various neural network topologies and configurations onto the hardware through adjustable parameters such as the number of neurons, the number of axons per core, and the precision of synaptic weights. This feature clearly distinguishes RANC\cite{RANC} from fixed-function neuromorphic systems implemented as ASIC chips, making it an ideal platform for research, prototyping, and optimization of novel SNN models before deployment on silicon.

Overall, the combination of distributed core architecture, event-driven communication, and high reconfigurability enables RANC\cite{RANC} to serve as an effective platform for implementing energy-efficient spiking neural networks, particularly in edge computing environments.

\subsubsection{Neuron and Synapse Mapping}
In the RANC\cite{RANC} ecosystem, the mapping of neurons and synapses is implemented based on a flexible crossbar-based architecture. Each Neuron Block emulates a crossbar connecting $N(a)$ input axons to $N(n)$ output neurons, allowing for a maximum of $N(a)$×$N(n)$ synaptic connections per core. All configuration parameters—including weights, activation thresholds, and voltage potentials—are stored centrally within the local core SRAM. The primary advantage of RANC\cite{RANC} over fixed-function architectures, such as IBM's TrueNorth, lies in its ability to customize the core structure to optimize mapping efficiency for specific applications.
\begin{itemize}
    \item \textbf{Vector-Matrix Multiplication (VMM)}: By supporting symmetric threshold behavior, RANC\cite{RANC} enables the mapping of signed calculations without the need for complex and resource-heavy feedback systems.This optimization reduces the required neuron and axon footprint by approximately \textbf{55-57\%}.
    \item \textbf{Convolutional Neural Networks (CNNs)}: RANC\cite{RANC} allows for the deployment of heterogeneous, large-scale cores (e.g., 1024 axons × 256 neurons). This mechanism significantly improves neuron utilization—increasing it from 0.7\% to 94.5\%—by minimizing the need for extensive input data repetition across multiple cores.
\end{itemize}

Through the optimization of core sizes to fully utilize BRAM primitives on the Alveo U250 FPGA, the architecture can scale to emulate as many as 259,072 distinct neurons and over 73.3 million distinct synapses on a single hardware device.

\subsubsection{Reconfigurability and Flexibility}

\begin{figure}[tb]
    \centering
    \includegraphics[width=0.7\linewidth]{Figure/RANC-Ecosystem.pdf}
    \caption{ Training, Testing, and Emulation environment of offered by RANC.}
    \vspace{-12pt}
    \label{fig:Training-Testing-and-Emulation-environment-of-offered-by-RANC}
\end{figure}

A key strength of the RANC\cite{RANC} architecture lies in its high degree of reconfigurability and flexibility, which enables efficient support for a wide range of Spiking Neural Network (SNN) models. Unlike fixed-function neuromorphic systems, RANC\cite{RANC} allows users to adapt both the network structure and hardware configuration according to specific application requirements.

As illustrated in Fig.~\ref{fig:Training-Testing-and-Emulation-environment-of-offered-by-RANC}, the RANC\cite{RANC} ecosystem provides a unified workflow that integrates training, simulation, and hardware emulation. Users can define multiple parameters, including network topology, number of neurons per core, number of inputs (axons), number of output classes, and the spike encoding scheme. These parameters directly influence how the SNN model is mapped onto the underlying hardware architecture, enabling flexible exploration of different design configurations.

At the architectural level, this flexibility is achieved through configurable neuron cores and programmable routing mechanisms. Each neuron core can be customized in terms of its internal structure and resource allocation, allowing the system to accommodate networks of varying sizes and complexities. In addition, the packet-based communication model enables dynamic routing of spike events between cores, removing the need for fixed interconnections and supporting diverse connectivity patterns.

Furthermore, the RANC\cite{RANC} framework supports multiple execution environments, as shown in Fig.~\ref{fig:Training-Testing-and-Emulation-environment-of-offered-by-RANC}, including a simulation environment, hardware emulation, and deployment on FPGA. This multi-stage pipeline allows models to be trained and validated in software before being translated into hardware configurations. The separation between training and deployment ensures that different system configurations can be evaluated efficiently without modifying the core hardware design.

Another important aspect of flexibility is the ability to reuse the same hardware platform for different SNN models. By simply modifying configuration parameters and input encoding schemes, researchers can experiment with various architectures without redesigning the system. This makes RANC\cite{RANC} particularly suitable for rapid prototyping and iterative development.

Overall, the reconfigurability and flexibility of RANC\cite{RANC} not only improve hardware utilization but also provide a practical framework for bridging the gap between SNN model development and hardware deployment. This characteristic is highly relevant to this thesis, where different SNN architectures are explored under a unified, hardware-aware design perspective.

\subsubsection{FPGA Implementation}
Field-Programmable Gate Arrays (FPGAs) provide a practical and flexible platform for implementing neuromorphic architectures such as RANC. Unlike application-specific integrated circuits (ASICs), FPGAs allow rapid prototyping and iterative design, making them particularly suitable for research and development of Spiking Neural Networks (SNNs). The reconfigurable nature of FPGAs aligns well with the design philosophy of RANC, enabling efficient mapping of neuron cores, synaptic connections, and routing networks onto hardware.

In the context of RANC, the architecture can be synthesized and deployed on FPGA devices by implementing neuron cores as parallel processing modules and the on-chip communication network as configurable routing logic. The use of packet-based dimension-order (XY) routing simplifies the hardware implementation, as spike events can be efficiently encoded, transmitted, and scheduled using digital logic. Furthermore, the global tick mechanism ensures synchronized operation across all processing elements, maintaining timing accuracy throughout the network. 
One of the key advantages of FPGA-based implementation is the ability to evaluate energy efficiency and performance trade-offs under realistic hardware constraints. By exploiting parallelism and event-driven computation, SNNs deployed on FPGA platforms can achieve significant reductions in power consumption compared to conventional GPU-based implementations. This makes FPGA-based neuromorphic systems particularly attractive for edge computing applications, where energy and latency are critical considerations.
However, FPGA implementations also introduce certain challenges. Resource limitations, such as the availability of Block RAM (BRAM) and Look-Up Tables (LUTs), can restrict the scale of deployable SNN models. Additionally, efficient mapping from high-level SNN models to hardware requires careful design to balance accuracy, latency, and hardware utilization.
In this thesis, although the proposed SNN model is developed with consideration for the RANC\cite{RANC} architecture, the actual deployment on FPGA hardware is beyond the scope of this work. Nevertheless, the compatibility of the proposed model with RANC\cite{RANC} highlights its potential for future implementation on reconfigurable hardware platforms, paving the way for energy-efficient and deployable fingerprint recognition systems.

In this thesis, although the proposed SNN model is developed with consideration for the RANC\cite{RANC} architecture, the actual deployment on FPGA hardware is beyond the scope of this work. Nevertheless, the compatibility of the proposed model with RANC\cite{RANC} highlights its potential for future implementation on reconfigurable hardware platforms, paving the way for energy-efficient and deployable fingerprint recognition systems.

\subsection{Vision Transformer \& SNN-ViT}
\subsubsection{Vision Transformer}
Vision Transformers (ViTs) have recently emerged as a powerful alternative to convolutional neural networks (CNNs) for image recognition by leveraging self-attention mechanisms to model global dependencies across an entire image. Unlike CNNs, which rely on local receptive fields and strong inductive biases such as translation invariance, ViTs enable direct modeling of long-range relationships between image regions, making them particularly effective for capturing distributed and non-local patterns. This characteristic is especially relevant for challenging data such as low-resolution and noisy capacitive images, where discriminative features may not be confined to local neighborhoods. However, ViTs typically require large-scale training data and incur high computational cost, which limits their applicability in resource-constrained environments. These limitations motivate further exploration of efficient variants, including the integration of ViTs with energy-efficient paradigms such as Spiking Neural Networks (SNNs) \cite{dosovitskiy2021imageworth16x16words}, \cite{touvron2021trainingdataefficientimagetransformers}.

The core idea of a Vision Transformer (ViT) is to treat an image as a sequence of patches, analogous to tokens in natural language processing. Given an input image $x \in \mathbb{R}^{H \times W \times C}$, the image is divided into a grid of non-overlapping patches, each of size $P \times P$. The total number of patches is $N = \frac{HW}{P^2}$. Each patch is then flattened into a vector and projected into a latent embedding space through a learnable linear transformation.

To retain spatial information, positional embeddings are added to the patch embeddings. In addition, a learnable classification token, denoted as $x_{\text{cls}}$, is prepended to the sequence. The final input sequence to the Transformer encoder can be written as:

\begin{equation}
Z_0 = [x_{\text{cls}}; x_1 E; x_2 E; \dots; x_N E] + E_{\text{pos}},
\end{equation}

where $x_i$ denotes the $i$-th image patch, $E$ is the patch embedding projection matrix, and $E_{\text{pos}}$ represents the positional embeddings. The resulting sequence of embeddings serves as the input to a standard Transformer encoder.

The Transformer encoder consists of a stack of identical layers, each composed of a multi-head self-attention (MHSA) mechanism and a position-wise feed-forward network (FFN), combined with residual connections and layer normalization. The self-attention mechanism allows each patch to attend to all other patches in the image, enabling the model to learn global contextual relationships. Specifically, attention weights are computed based on the similarity between query, key, and value representations derived from the input embeddings, allowing the model to dynamically focus on relevant regions of the image. The use of multiple attention heads further enhances the model’s capacity to capture diverse patterns and interactions across different representation subspaces.

\begin{figure}[tb]
    \centering
    \includegraphics[width=0.95\linewidth]{Figure/model_scheme.pdf}
    \caption{Vision Transformer Overview.}
    \vspace{-12pt}
    \label{fig:ViT_scheme}
\end{figure}

After passing through the Transformer encoder, the representation corresponding to the classification token is used as a global image descriptor. This representation is then fed into a multi-layer perceptron (MLP) head to produce the final class prediction. Compared to CNN-based approaches, which progressively aggregate local features through convolution and pooling operations, ViTs perform global reasoning at every layer, providing a more flexible and expressive modeling framework.

Despite these advantages, Vision Transformers also present several limitations. In particular, the lack of strong inductive biases makes them less data-efficient than CNNs, often requiring large-scale datasets to achieve competitive performance. Additionally, the quadratic complexity of the self-attention mechanism with respect to the number of patches leads to increased computational and memory costs, which can be prohibitive for deployment on edge devices. These challenges are especially relevant in the context of this thesis, where energy efficiency and hardware constraints are key considerations.

\begin{figure}[tb]
    \centering
    \includegraphics[width=0.3\linewidth]{Figure/20201002_selected_attention_examples.pdf}
    \caption{Representative examples of attention from the output token to the input space.}
    \vspace{-12pt}
    \label{fig:Representative_examples_of_attention_from_the_output_token_to_the_input_space}
\end{figure}

In the context of capacitive fingerprint recognition, the ability of ViTs to model global relationships is particularly beneficial. Capacitive images are characterized by low resolution, high noise levels, and significant variability in touch position and pressure, making it difficult for purely local feature extractors to capture discriminative patterns. By leveraging self-attention, ViTs can integrate information across the entire spatial domain, potentially improving robustness and classification accuracy. Nevertheless, the high energy consumption of conventional ViTs motivates the exploration of more efficient architectures. This leads to the investigation of Spiking Neural Network-based Vision Transformers (SNN-ViT), which aim to combine the representational power of Transformers with the energy efficiency of event-driven computation.

\subsubsection{SNN-ViT}
Recent advances in neuromorphic computing have motivated the integration of Transformer architectures into the spiking domain, leading to the emergence of Spiking Vision Transformers (SNN-ViTs). These models aim to combine the global representation capability of Vision Transformers (ViTs) with the energy efficiency of Spiking Neural Networks (SNNs).

Unlike conventional ViTs, which operate on continuous-valued activations, SNN-ViTs process information through discrete spike events over time. This requires redesigning key Transformer components—including patch embedding, self-attention, and feed-forward layers—to be compatible with spike-based computation.

Several approaches have been proposed in the literature to bridge the gap between Vision Transformers and Spiking Neural Networks. One early direction focuses on ANN-to-SNN conversion for Transformer architectures, where a pre-trained Vision Transformer is transformed into a spiking equivalent by mapping continuous activations to firing rates. While this approach benefits from leveraging well-optimized ANN models, it inherits several limitations commonly observed in CNN-to-SNN conversion \cite{diehl2015unsupervised, pfeiffer2018deep}. In particular, the reliance on rate-based encoding often leads to long inference latency due to the need for multiple simulation time steps. Moreover, the complex structure of self-attention mechanisms introduces additional challenges, resulting in potential accuracy degradation when approximated in the spiking domain.

To overcome these limitations, more recent studies have shifted toward directly trained spiking Transformers, where models are trained end-to-end using surrogate gradient methods \cite{neftci2019surrogate}. This paradigm enables the network to learn spike-based representations natively, without relying on intermediate ANN formulations. A representative example is SpikeFormer, which introduces a fully spike-driven Transformer architecture by replacing conventional activation functions with spiking neurons while preserving the overall Transformer structure \cite{zhou2023spikeformer}. Such approaches have demonstrated competitive performance on standard benchmarks, while significantly improving energy efficiency due to their event-driven computation.

In addition to these methods, hybrid architectures have been proposed to further balance efficiency and representational power. For instance, models such as SpikeAtConv \cite{SpikeAtConv} integrate convolutional layers with spiking attention mechanisms, allowing the network to capture both local spatial features and global contextual relationships. By combining the strengths of convolution and attention, these hybrid designs aim to reduce the computational overhead associated with pure Transformer-based models while maintaining strong feature extraction capability.

Despite these advancements, existing SNN-ViT approaches are predominantly evaluated on high-resolution datasets such as ImageNet \cite{ImageNet}, where abundant spatial information is available. Consequently, their effectiveness on low-resolution and noisy data—such as capacitive fingerprint images—remains largely unexplored. This limitation highlights an important research gap in the current literature.

\begin{figure}[tb]
    \centering
    \includegraphics[width=1\linewidth]{Figure/ssa.pdf}
    \caption{ Illustration of vanilla self-attention (VSA)\cite{dosovitskiy2021imageworth16x16words} and Spiking Self Attention (SSA)\cite{zhou2023spikeformer}}
    \vspace{-12pt}
    \label{fig:Spiking-Self-Attention}
\end{figure}

To address this gap, this thesis adopts a SpikeFormer-based Spiking Vision Transformer as the core architecture for the CapFingerId task. Specifically, the proposed model largely follows the original SpikeFormer design \cite{zhou2023spikeformer}, while introducing task-specific adaptations to better suit low-resolution capacitive data. In particular, the number of Spiking Transformer encoder blocks and the feature embedding dimensions are carefully adjusted to balance representational capacity and computational efficiency under resource-constrained settings.

By retaining the core spike-driven self-attention mechanism and direct training strategy with surrogate gradients, the model preserves the advantages of SpikeFormer in capturing global dependencies. At the same time, the architectural scaling enables more efficient processing of small, noisy inputs, making it more suitable for capacitive fingerprint classification. This design choice allows the thesis to focus on evaluating the applicability and effectiveness of Spiking Vision Transformers in a new problem domain, rather than proposing a fundamentally new architecture.

\subsubsection{Energy Consumption Estimation}
\label{sec:energy_estimation}

Following the methodology adopted in recent Spiking Transformer \cite{zhou2023spikeformer} studies, the
energy consumption of the proposed model is estimated analytically rather than
measured directly on conventional GPU hardware. This estimation is based on the
computational cost of individual operations, where a multiply--accumulate (MAC)
operation consumes significantly more energy than an accumulate (AC) operation.
According to the energy model proposed by Horowitz \cite{horowitz20141} for a 45 nm CMOS process, the
energy cost of a single MAC and AC operation is assumed to be
$E_{\mathrm{MAC}} = 4.6~\mathrm{pJ}$ and
$E_{\mathrm{AC}} = 0.9~\mathrm{pJ}$, respectively.

For an artificial neural network (ANN), the energy consumption of the
$l^{\mathrm{th}}$ layer is estimated as

\begin{equation}
E_{\mathrm{ANN}}^{(l)}
=
N_{\mathrm{MAC}}^{(l)}
E_{\mathrm{MAC}},
\end{equation}

where $N_{\mathrm{MAC}}^{(l)}$ denotes the number of multiply--accumulate
operations required by the layer. Consequently, the total energy consumption of
an ANN is

\begin{equation}
E_{\mathrm{ANN}}
=
\sum_{l=1}^{L}
N_{\mathrm{MAC}}^{(l)}
E_{\mathrm{MAC}},
\end{equation}

where $L$ is the total number of layers.

For a spiking neural network (SNN), only active spikes trigger accumulation
operations. Therefore, the computational cost depends on the firing rate of each
layer. The firing rate of layer $l$ is defined as

\begin{equation}
r^{(l)}
=
\frac{N_{\mathrm{spike}}^{(l)}}
{N_{\mathrm{neuron}}^{(l)}},
\end{equation}

where $N_{\mathrm{spike}}^{(l)}$ is the total number of emitted spikes and
$N_{\mathrm{neuron}}^{(l)}$ is the total number of neurons evaluated over all
simulation time steps.

The number of synaptic accumulation operations (SOPs) is then estimated by

\begin{equation}
N_{\mathrm{SOP}}^{(l)}
=
r^{(l)}
T
N_{\mathrm{MAC}}^{(l)},
\end{equation}

where $T$ denotes the number of simulation time steps.

The energy consumption of a spiking layer is therefore

\begin{equation}
E_{\mathrm{SNN}}^{(l)}
=
N_{\mathrm{SOP}}^{(l)}
E_{\mathrm{AC}}.
\end{equation}

For a general spiking neural network, some layers may operate on continuous-valued
inputs and therefore require MAC operations, while the remaining layers process
spike trains and only perform AC operations. Let $\mathcal{M}$ denote the set of
layers requiring MAC operations and $\mathcal{A}$ denote the set of layers
performing AC operations. The total energy consumption of an SNN can then be
estimated as

\begin{equation}
E_{\mathrm{SNN}}
=
\sum_{l\in\mathcal{M}}
N_{\mathrm{MAC}}^{(l)}
E_{\mathrm{MAC}}
+
\sum_{l\in\mathcal{A}}
N_{\mathrm{SOP}}^{(l)}
E_{\mathrm{AC}}.
\end{equation}

Substituting Eq.~(5) into the above expression yields

\begin{equation}
E_{\mathrm{SNN}}
=
\sum_{l\in\mathcal{M}}
N_{\mathrm{MAC}}^{(l)}
E_{\mathrm{MAC}}
+
E_{\mathrm{AC}}
\sum_{l\in\mathcal{A}}
r^{(l)}
T
N_{\mathrm{MAC}}^{(l)}.
\end{equation}

In many SNN architectures, such as SpikeFormer~\cite{zhou2023spikeformer}, only
the input encoding layer operates on continuous-valued inputs, whereas all
subsequent layers process spike trains. In this case,
$\mathcal{M}$ contains only the first layer and Eq.~(8) naturally reduces to the
energy estimation formula commonly adopted in the SpikeFormer \cite{zhou2023spikeformer} literature.

\end{document}